\documentclass[a4paper]{article}

% Basic packages
\usepackage[fontset=fandol]{ctex}
\usepackage{amsmath, amssymb}
\usepackage{graphicx}
\usepackage[margin=2.5cm]{geometry}
\usepackage[most]{tcolorbox}
\usepackage{etoolbox}
\usepackage{listings}
\usepackage{hyperref}
\usepackage{booktabs}
\usepackage{subcaption}
\usepackage{float}
\usepackage{tikz}
\IfFileExists{pgfplots.sty}{
    \usepackage{pgfplots}
    \pgfplotsset{compat=1.18}
}{}

% Highlight boxes
\newtcolorbox{knowledgebox}[1]{
    enhanced, colback=blue!5!white, colframe=blue!75!black, colbacktitle=blue!75!black,
    coltitle=white, fonttitle=\bfseries, title=#1, attach boxed title to top left={yshift=-2mm, xshift=2mm},
    boxrule=1pt, sharp corners
}

\newtcolorbox{importantbox}[1]{
    enhanced, colback=yellow!10!white, colframe=yellow!80!black, colbacktitle=yellow!80!black,
    coltitle=black, fonttitle=\bfseries, title=#1, sharp corners
}

\newtcolorbox{warningbox}[1]{
    enhanced, colback=red!5!white, colframe=red!75!black, colbacktitle=red!75!black,
    coltitle=white, fonttitle=\bfseries, title=#1, sharp corners
}

\newtcolorbox{dialoguebox}[1]{
    enhanced, breakable, colback=green!4!white, colframe=green!45!black, colbacktitle=green!45!black,
    coltitle=white, fonttitle=\bfseries, title=#1, sharp corners
}

% Listings
\lstset{
    language=Python,
    basicstyle=\ttfamily\small,
    keywordstyle=\color{blue},
    stringstyle=\color{red!60!black},
    commentstyle=\color{green!60!black},
    breaklines=true,
    frame=single,
    numbers=left,
    numberstyle=\tiny\color{gray},
    captionpos=b,
    extendedchars=false
}

% Front-page metadata
\newcommand{\notetitle}{课程笔记：语言模型智能体的工具使用}
\newcommand{\noteauthors}{Bowen-mantou（整理）}
\newcommand{\notedate}{2026-10-23}
\newcommand{\videochannel}{B 站「常青藤名校课」· CMU 11-768 AI Agents（英文原声字幕版）}
\newcommand{\videopublishdate}{2026-09-09}
\newcommand{\videoduration}{57:22}
\newcommand{\videourl}{https://www.bilibili.com/video/BV1a6Yx62EH4?p=16}
\newcommand{\videocoverpath}{cover.jpg}

\begin{document}

\begin{titlepage}
\centering
\vspace{1.2cm}
{\huge\bfseries \notetitle\par}
\vspace{0.8cm}
{\large \noteauthors\par}
\vspace{0.3cm}
{\large \notedate\par}
\vspace{1.2cm}

\ifdefempty{\videocoverpath}{
{\small 视频封面因版权未收录，见下方视频链接。\par}
}{
\par
}

\vfill
\begin{tcolorbox}[width=0.9\textwidth, colback=black!2!white, colframe=black!60, sharp corners]
\textbf{视频作者/频道}：\videochannel\par
\textbf{发布日期}：\videopublishdate\par
\textbf{视频时长}：\videoduration\par
\textbf{视频链接}：
\ifdefempty{\videourl}{
未填写
}{
\href{\videourl}{\nolinkurl{\videourl}}
}
\par\textbf{PDF制作Skill}：\href{https://github.com/wdkns/wdkns-skills}{\nolinkurl{https://github.com/wdkns/wdkns-skills}}
\end{tcolorbox}
\end{titlepage}

\tableofcontents
\newpage

%% --- 正文内容开始 --- %%

\section{引言：工具是智能体与语言模型的分水岭}

语言模型本质上是一台「预测下一个 token」的机器：它的参数在训练结束后就冻结了，既看不到今天的信息，也无法在真实世界里动手。工具（tool）的意义，就是给这台机器开一个通往外部程序的接口——模型输出一段结构化的调用请求，外部程序真正执行，再把结果作为观察（observation）回填到上下文里。主讲人把这一点讲得很直接：\textbf{工具使用是智能体（agent）与语言模型最根本的区别}，没有工具调用，就只剩下「一次性的文本问答」。

一个容易被忽略的细节是：\textbf{工具调用的形式是标准化的，但工具下游能做什么完全由你决定}。搜索、计算器、浏览器、Python 解释器、任意第三方 API，甚至调用一个电商接口替用户下单，都共用同一套「模型发出调用、系统执行、结果回填」的协议。标准化发生在协议层，能力和风险则发生在实现层。

\begin{knowledgebox}{本讲路线图}
为什么需要工具（两类收益）$\rightarrow$ 工具的分类与形态 $\rightarrow$ 函数调用、代码调用、并行调用三种范式 $\rightarrow$ 工具调用在模型内部如何用 token 表达与训练 $\rightarrow$ 解析、校验、分派、结果配对 $\rightarrow$ 受限解码（让非法调用不可能生成）$\rightarrow$ REST 与 MCP 两种接口工程 $\rightarrow$ 工具使用如何评测。
\end{knowledgebox}

\subsection{本章小结}
工具是智能体接入外部世界的唯一接口；本讲从「为什么需要」讲到「工程上怎么接、怎么保证不出错、怎么评测」。理解这一点，是理解后续长上下文、记忆、规划等所有 agent 能力的前提。

\section{工具的两大收益：扩展能力与便利化}

主讲人引用了课程阅读材料（Neubig、Fried 与 Zora Wang 合著的工具使用综述）中的分类，把工具的价值分成两类，这个二分法值得记住。

\textbf{第一类是「扩展」（extending）}：语言模型在原理上做不到的事情，靠工具能做到。最典型的是时效性——模型的参数停留在训练截止时刻，问它「四个月前的今天发生了什么」它只能编；而调用一个外部时间/检索函数就能拿到真实答案。同一类还包括访问私有状态（用户资料、日历、订单）以及对外部环境做出改变（浏览网页、调用某个 API 下单）。

\textbf{第二类是「便利化」（facilitating）}：语言模型其实能做，但让工具做更快、更可靠、更便宜。例如多位数乘法，模型要靠长链推理一步步算，速度慢且可能算错，而计算器可以瞬间给出精确结果。一句话概括：\textbf{扩展是「不能变成能」，便利化是「能变成又快又准」}。主讲人在课上给出的判据是：只有当工具的收益超过它的延迟、成本、失败风险时，才应该调用工具。

\subsection{工具的分类：从文本回复到自定义 API}

课堂上把常见工具归成了几大类，几乎任何新工具都能放进其中某一类：

\begin{enumerate}
\item \textbf{文本回复（textual responses）}：这是最「无聊」的一类，模型直接给出答案。实现上有两种方式——提供 \texttt{finish(text)} 这样的结束工具，或者干脆让模型不调用任何工具、结束 agent 循环。
\item \textbf{信息检索（web search / RAG）}：检索外部信息再基于其生成答案。主讲人在这里纠正了一个流行论调：所谓「RAG 已死」并不成立，\textbf{RAG 恰恰是因为无处不在才让人觉得它「死了」}——检索增强已经是几乎所有应用的默认形态。
\item \textbf{代码执行（code execution）}：把一段代码交给解释器去算，适合精确计算和数据处理。
\item \textbf{图像生成（image generation）}：与文本类工具不同，它的接口是「图像」。一个研究味很浓的细节是：模型往往会在调用前把用户 query 大幅改写/扩展，甚至下载图片时你会看到与用户原话完全不同的 caption——\textbf{改写 query 本身就是工具调用能力的一部分}。
\item \textbf{自定义函数（custom functions）}：如「把香蕉、牛奶、咖啡加进购物车」这类业务 API，本质上是把某个垂直能力包装成一个函数。
\end{enumerate}

课堂用 ChatGPT 的演化说明这件事有多普遍：它最初只是「文本进、文本出」，如今已经是一个「everything app」——写代码、渲染 3D、生成图像、生成 PDF、连接 Slack 等第三方服务。用户感知到的是功能列表，工程上则是背后挂载的一长串工具与 API。

\subsection{本章小结}
工具分两类价值：扩展模型做不到的能力（时效、私有状态、真实动作），以及替代模型做得不够好的能力（精确计算）。工具形态可以归为文本回复、检索、代码执行、图像生成、自定义 API 五类；「RAG 无处不在」和「query 改写」是两个值得记住的细节。

\section{三种工具调用范式：函数、代码与并行}

\subsection{函数调用：给模型一张 API 清单}

最直接的范式是：把 50--70 个 API 函数连同各自的 schema 一起交给模型，模型每一步选一个（或并行选多个）函数调用。这是当前工业界最主流的做法，优点是权限边界清晰、每个调用都能审阅；缺点是当任务需要「查天气 $\rightarrow$ 查价格 $\rightarrow$ 换算汇率 $\rightarrow$ 再查价格」这种多步组合时，交互轮数会变得很多，延迟和成本随之上升。

\subsection{代码即元工具：CodeAct 的思路}

第二种范式把代码本身当作工具。代码是极其表达力强的媒介：一段程序可以看作一棵函数调用树，还带循环、条件与变量；调用 \texttt{pandas} 就等于一次性打开了整个库的工具箱。因此让 agent 直接「写一段程序调用工具」，往往比逐个 API 调用更自然。

主讲人举了一个经典例子：判断在美国、日本、德国、印度买某款手机哪个最便宜。按传统方式，模型要一次一次调用「查汇率」「查手机价格」「算税」「换算」等工具，来回很多轮；而写成一段程序，一次生成就能把整套流程表达完——这是 CodeAct 的核心动机。论文报告的结果也很直接：\textbf{任务成功率上升、平均交互轮数下降}，而且这个收益不限于「本来就需要写代码」的任务。

\subsection{代码调用的代价：为什么不能一切皆代码}

代码能力越强，风险面越大。课堂讨论了四类问题：

\begin{enumerate}
\item \textbf{死循环}：模型可能写出无限循环，系统必须有超时/中断机制兜底。
\item \textbf{供应链安全}：模型可能引入被投毒的第三方库，一旦中招，整个 agent 系统会被攻破——这是当前 agent 安全里非常现实的一类风险。
\item \textbf{难约束}：动作空间大、后果不可预测，很难像 API 那样做细粒度权限校验。
\item \textbf{资源影响}：代码可以吃掉大量内存、磁盘和 CPU，需要显式限制资源。
\end{enumerate}

因此，一旦把工具调用升级为「程序化调用」，就必须配套沙箱（sandbox）与资源限额。课程把它留到安全专题展开，但结论先给出来了：\textbf{能力越强，容器就要越硬}。

\subsection{并行工具调用：一次生成，多个动作}

当多个工具调用彼此独立时（例如同时查天气、查日历、查航班状态），可以让模型在单次响应里发出多个工具调用块，系统用 \texttt{asyncio.gather} 一类原语并发执行，再把多个结果一起回填。反之，有依赖关系的调用（先查用户 ID，再用 ID 查订单，再退款）必须串行——\textbf{能否并行取决于调用之间是否存在数据依赖}。

一个反直觉的现象是：\textbf{按任务完成成本算，更贵的模型反而可能更便宜}。原因有二：更强的模型一次选对方案、少走弯路；更强的模型并行工具调用能力更好，能一次读多个文件、同时看多张图。主讲人指出，并行调用的能力在最近半年才明显变好，一个重要推动力是强化学习——训练时对完成时长/冗长程度施加惩罚，会天然激励模型把独立调用并行化。

\subsection{本章小结}
函数调用胜在可控，代码调用胜在表达力与轮数，并行调用胜在延迟。三者的选择标准是「权限、依赖、时延」三件事：权限敏感用函数，多步组合用代码，相互独立则并行；而代码调用必须配沙箱。

\section{工具调用在模型内部：token、协议与模板}

\subsection{工具调用是 token 序列}

工具调用不是「另一套 API」，而是模型生成的 token 序列。普通文本续写是「今天天气是晴天」；工具续写是 \texttt{<tool\_call>get\_weather(...)</tool\_call>} 这样的结构；实际输出里通常两者兼有——先对用户说一句话，再发出工具调用。这正是 ReAct 循环在今天的落地方式：\textbf{ReAct 不再是写在外部代码里的 while 循环，而是一次 completion 里同时包含思考、发言与调用}。对带长思考的推理模型，输出还可能分成「隐式的 thinking token」「给用户看的 message」「真正的 tool call」三段。

\subsection{训练视角：格式 token 从哪来}

课堂提问：「这些特殊 token 需要在预训练时就进词表吗？」主讲人的回答是：不一定。更常见的做法是在\textbf{中训练（mid-training）}阶段引入工具调用格式，然后才做强化学习；预训练在整张互联网上做时，词表里未必有这些 token。这一点对做后训练的人很重要：\textbf{工具调用格式是一项需要专门训练的能力，而不是凭空泛化出来的}。

\subsection{每个模型的协议都不一样}

同样的工具调用，不同模型的编码格式不同：Qwen 用 ChatML 风格的特殊标记；Mistral 有自己的一套 tool call 标记；DeepSeek 把它叫 function call，用自研的 DSML（DeepSeek Markup Language）标记；Anthropic 又是另一种格式。工程结论有两条：

\begin{enumerate}
\item \textbf{微调时必须匹配基座模型的原生格式}——不要试图让模型学一个它没见过的格式；
\item \textbf{解析代码不能跨模型复用}——为 DeepSeek 写的解析器未必能解析 Qwen 的输出，反之亦然。
\end{enumerate}

好消息是这类模板已经基本标准化（Hugging Face 的 \texttt{apply\_chat\_template} 会按模型注入正确格式），但知道底层发生了什么，能避免很多「看起来莫名奇妙」的 bug。

\subsection{本章小结}
工具调用是在 token 层面生成的结构化文本，格式 token 通常在中训练阶段引入；不同模型协议不同，微调必须匹配原生模板，解析器不能跨模型照搬。

\section{从模型输出到真正执行：解析、分派与结果配对}

模型吐出调用之后，agent 框架要完成一条固定的流水线：

\begin{enumerate}
\item \textbf{解析（parse）}：把模型输出解析成「工具名 + 参数」；
\item \textbf{注册与解析（register / resolve）}：维护一张「工具名 $\rightarrow$ 可执行函数」的映射表，框架通常提供注册新工具的接口；
\item \textbf{校验（validate）}：检查工具是否存在、参数是否符合 schema；
\item \textbf{执行（execute）}：调用真实函数，捕获异常；
\item \textbf{返回观察（observation）}：把结果包装成一条带调用 ID 的消息，回填进对话历史。
\end{enumerate}

这里有两类失败要区分：\textbf{校验失败}是模型选错了工具或参数不合法；\textbf{执行失败}是参数看起来合法、但运行时出错——典型例子是「执行 Python」工具，任何字符串都能通过校验，但真跑起来可能报语法错误或异常。工程上要把两者分开处理，前者可以提示模型改正，后者更适合把异常信息作为观察回填。

另一个容易踩的坑是 \textbf{调用 ID 与结果 ID 的配对}。每次工具调用会被分配一个 ID，对应的结果必须带同一个 ID；并行调用时靠它区分「哪个结果对应哪个调用」。多数模型对此宽容，但 Anthropic 的接口在遇到「有调用、无结果」的历史时会直接拒绝继续生成。这对应一类现实故障：agent 发出调用后程序崩溃、结果没写回，恢复会话时历史不完整——修复方式是在恢复时补齐缺失的结果（哪怕是错误信息），而不是把残缺历史直接发给模型。

\subsection{本章小结}
工具调用是一条「解析 $\rightarrow$ 校验 $\rightarrow$ 分派 $\rightarrow$ 执行 $\rightarrow$ 回填」的流水线；要区分校验失败与执行失败，并保证每个调用都有 ID 匹配的结果，否则会制造出模型无法继续的残缺历史。

\section{受限解码：让不合法的工具调用不可能生成}

\subsection{问题：模型不保证输出合法 JSON}

语言模型不保证生成格式正确的工具调用——缺一个右括号的 JSON 就是典型例子。很多人第一反应是「写个脚本把括号补上」，但这既不可靠也不优雅。正确思路是把约束前置到解码阶段。

工具调用的合法性可以拆成三层约束：

\begin{enumerate}
\item \textbf{语法（syntax）}：必须是合法 JSON；
\item \textbf{形状（shape）}：必填字段齐全；
\item \textbf{类型与单位（types / units）}：例如 \texttt{city} 必须是字符串，\texttt{units} 只能是摄氏度或华氏度。这三层都可以用 JSON schema 表达。
\end{enumerate}

\subsection{形式语言视角：从正则文法等价到下推自动机}

课堂用形式语言理论解释了为什么这件事可以做到「生成即合法」。正则表达式对应\textbf{正则文法}，可以被\textbf{有限状态自动机（FSA）}识别：每读入一个 token，状态转移一次。但 JSON schema 不是正则语言——JSON 有嵌套括号、逗号分隔的成员，存在递归结构，需要用\textbf{上下文无关文法（CFG）}描述；识别 CFG 需要比有限状态自动机更强的机器：\textbf{下推自动机（PDA）}，也就是在状态机上再加一个栈，用入栈/出栈记住当前嵌套深度。

\subsection{xGrammar：在 logits 上做掩码}

工程上最流行的实现是 \textbf{xGrammar}（CMU 机器学习系的工作，已被广泛集成进各家推理库）。其核心机制是：每生成一个 token 之前，查当前的语法状态，允许的 token 保留原 logits，不允许的 token 置为负无穷，再重新归一化：

$$
p'(x_t = v \mid x_{<t}) =
\frac{\exp(z_v)\cdot \mathbb{1}[v \in \mathcal{V}_{\text{valid}}(s_t)]}
{\sum_{v'}\exp(z_{v'})\cdot \mathbb{1}[v' \in \mathcal{V}_{\text{valid}}(s_t)]},
$$
其中 $z_v$ 是模型对 token $v$ 的原始 logits，$s_t$ 是生成到第 $t$ 步时的语法状态（包含下推自动机的栈），$\mathcal{V}_{\text{valid}}(s_t)$ 是在该状态下合法的 token 集合，$\mathbb{1}[\cdot]$ 是指示函数。这样生成的序列在语法上永远合法。

性能上，xGrammar 做了两类优化：先预计算「永远合法」和「永远非法」的词表项，只有随上下文变化的 token 才在运行时按栈状态计算；同时把 CPU 上的语法计算与 GPU 上的推理重叠起来，从而把约束解码的额外开销压得很低。

\subsection{唯一的漏洞：token 预算耗尽}

即使开启受限解码，仍有一种情况会生成不完整的 JSON：模型的输出 token 上限到了，而当前状态虽然仍在一条合法路径上，却还没走到接受状态——于是 JSON 被截断。理论上可以在生成时统计「还差多少 token 才能收尾」，但工业实现里并不总是做这件事。知道这个边界，面试和工程排查时都用得上。

\subsection{本章小结}
工具调用的合法性分语法、形状、类型三层，都能由 JSON schema 描述；JSON 需要上下文无关文法，识别它用下推自动机；xGrammar 通过对 logits 做掩码实现「生成即合法」，其唯一边界是 token 预算耗尽导致的截断。

\section{工具接口的工程化：REST API 与 MCP}

\subsection{三种把工具接进 agent 的方式}

\begin{enumerate}
\item \textbf{程序化调用}：给 agent 一个 Python 环境和一个库（如 \texttt{weatherlib.py}），让它写代码调用函数。简单直接，但不能远程调用，也不兼容标准工具调用通道。
\item \textbf{REST API}：用 FastAPI 之类的框架把函数暴露成 HTTP 接口，配一个 API key，写一个装饰器就把函数变成端点。好处是：可以用 key 限制访问；而且整个 API 会自动生成 JSON schema，可以直接交给模型当工具定义。coding agent 若能用 Bash，也可以直接用 \texttt{curl} 调 REST 接口。
\item \textbf{MCP（Model Context Protocol）}：由 Anthropic 提出的「给模型提供上下文/工具」的标准协议，架构是 host（AI 应用）连接多个 client，每个 client 再连一个 MCP server（本地 stdio 或远程 socket）。
\end{enumerate}

\subsection{MCP 真正的价值：多一层凭证隔离}

主讲人坦言自己最初不理解 MCP 的必要性：既然 REST 已经能调 API，为什么还要多一个「跑程序的协议」？答案在于\textbf{安全边界}。MCP 有两层 key：

\begin{itemize}
\item \textbf{MCP API key}：交给运行 agent 的进程，用来连 MCP server；
\item \textbf{上游 API key}（如 GitHub token）：只交给 MCP server，不交给 agent。
\end{itemize}

这样即使 agent 把 MCP key 泄露了，泄露的东西也基本无害；真正的上游凭证始终不在 agent 手里。反例很好记：如果把 GitHub token 直接给 agent，它可能某天「灵机一动」把 token 推到公开仓库，账号与所有仓库就都沦陷了。除了凭证隔离，MCP 与 REST 还在发现方式（OpenAPI 文档 vs 专门的 list tools 接口）、调用传输（HTTP vs stdio/socket）、能力范围（resources/prompts/extensions）上不同；生态上已有官方 MCP registry 和大量聚合站点，FastMCP 则让「把 REST 快速变成 MCP server」变得很容易。

\subsection{本章小结}
工具接入有三种方式：程序化调用最直接、REST 最通用、MCP 最利于权限隔离。MCP 的核心价值不是「多一层协议」，而是把上游凭证留在 server 侧、只把无害的代理 key 交给 agent。

\section{工具使用的评测：把「会调工具」变成可测量的指标}

工具调用是 agent 能力的前置条件，因此需要单独评测。最著名的是 \textbf{Berkeley Function Calling Leaderboard（BFCL）}，从 v1 到 v4 逐步覆盖：单轮工具调用 $\rightarrow$ 多轮工具调用 $\rightarrow$ agentic 工具调用（多步任务）$\rightarrow$ 鲁棒性（参数幻觉、格式敏感性等）。

评测时要分维度看「整条工具使用链路」：

\begin{enumerate}
\item \textbf{选对工具}：面对任务，选择的工具是否真的能完成它；
\item \textbf{参数正确}：参数值是否正确（包括单位、格式）；
\item \textbf{顺序正确}：多步调用里次序是否合理；
\item \textbf{端到端准确率}：任务最终是否完成；
\item \textbf{效率}：完成需要多少 token、多少时延、多少成本；
\item \textbf{可靠性}：超时是否重试、工具失败能否换一个可用工具；
\item \textbf{安全性}：面对对抗性输入，是否会被诱导做出不当调用。
\end{enumerate}

课堂上还展示了一个非常实用的对比：同一个模型（字幕称 GLM 5.3，具体型号以原视频为准），换推理供应商后工具调用错误率可以从约 15\% 落到 0.01--0.05\%。原因包括：不同供应商的解码/推理实现不同；量化等级不同（主讲人的经验是 FP8 明显好于 FP4，FP4 省成本但更像「另一个模型」）；投机解码是否无损；是否实现了约束解码（vLLM、SGLang 支持语法约束，有些供应商没用）；以及供应商本身的可靠性（服务挂掉时调用自然失败）。这条对工程选型意义很直接：\textbf{选模型不仅是选权重，也是选推理栈}。

\subsection{本章小结}
工具使用可以拆成七个可测维度（选择、参数、顺序、端到端、效率、可靠性、安全）；BFCL 是标准基准；同一模型在不同推理供应商上的工具调用可靠性差异巨大，量化、投机解码和约束解码实现是主因。

\section{总结与延伸}

本讲的收尾把工具使用概括成一个分层系统：工具为 agent 增加检索、执行、生成、应用等能力，是 agent 之所以为 agent 的关键；但这一切不是免费的——你需要把 schema 放进请求、让模型生成调用、解析并校验、分派执行、按 ID 配对结果，还可以用语法约束解码保证格式合法；而语法合法并不等于语义正确（一段能被解析的 Python 仍然可能算错）。接口层面有 REST 与 MCP 两条路，编排层面则要处理依赖、并发与评测。

\begin{importantbox}{把本讲压缩成五句话}
\begin{enumerate}
\item 工具让语言模型从「生成文本」变成「对外部世界采取行动」；收益分扩展能力与便利化两类。
\item 工具调用有三种范式：函数清单、代码即元工具（CodeAct）、并行调用；代码表达力最强但必须沙箱化。
\item 工具调用是 token 序列，格式由各模型自己的协议决定；微调要对齐原生模板。
\item 执行侧是「解析、校验、分派、执行、回填」的流水线，call ID 必须与结果配对。
\item 受限解码（如 xGrammar）能让非法 JSON 无法生成，但只是语法约束；接口用 REST 或 MCP，评测看 BFCL 与多维度指标。
\end{enumerate}
\end{importantbox}

课程预告下一讲进入「长上下文建模」：既然工具会不断往上下文里塞观察，如何管理这些越来越长的上下文就是下一个核心问题。作业方面，第一次作业要求结合本讲写一段反思，并在 Canvas 上提交。

\textbf{延伸思考}：(1) 从训练视角看，工具层其实就是 RL 的「环境」——工具 schema、错误信息、沙箱边界都会影响 rollout 的质量与可验证性；(2) 从评测视角看，BFCL 这类基准测的是「工具使用」而不是「完成任务」，两者之间还有很大 gap，agent 的端到端评测仍然更难；(3) 从工程视角看，MCP 的凭证隔离思路值得推广到所有「让模型接触真实系统」的场景——把危险能力留在服务端，只给模型最小权限。

\end{document}
